\documentclass{scrartcl}
\usepackage{graphicx}  % required for inserting images

% for cyrillic and Bulgarian language
\usepackage[utf8]{inputenc}
\usepackage[T2A]{fontenc}
\usepackage[bulgarian]{babel}

% math packages
\usepackage{amsmath}
\usepackage{float}
\usepackage{amsfonts}
\usepackage{amssymb}
\usepackage{amsthm}
\usepackage{cancel}

\usepackage{ragged2e}  % adds FlushLeft

\usepackage{listings}
\usepackage{xcolor}

\definecolor{codegreen}{rgb}{0,0.6,0}
\definecolor{codegray}{rgb}{0.5,0.5,0.5}
\definecolor{codepurple}{rgb}{0.58,0,0.82}
\definecolor{backcolour}{rgb}{0.95,0.95,0.92}

\lstdefinestyle{mystyle}{
    % caption=Example C++,
    label={lst:listing-cpp},
    language=C++,
    backgroundcolor=\color{backcolour},   
    commentstyle=\color{codegreen},
    keywordstyle=\color{magenta},
    numberstyle=\tiny\color{codegray},
    stringstyle=\color{codepurple},
    basicstyle=\ttfamily\footnotesize,
    breakatwhitespace=false,         
    breaklines=true,                 
    keepspaces=true,                 
    numbers=left,                    
    numbersep=5pt,                  
    showspaces=false,                
    showstringspaces=false,
    showtabs=false,                  
    tabsize=2,
}
\lstset{style=mystyle}

\usepackage[colorlinks=true, linkcolor=blue, urlcolor=cyan]{hyperref}

\newtheorem{theorem}{Теорема}
\newtheorem{definition}{Дефиниция}

\title{Процеси и комуникация между тях в операционната система}
\author{Ивайло Андреев + Gemini 3 Flash}
\date{25 май 2026}

\begin{document}
\maketitle

\tableofcontents

\newpage

\begin{FlushLeft}

\section{Процеси и комуникационни канали --- основни абстракции в ОС}

\subsection{Същност на Операционната система}
Операционната система (ОС) е системен софтуер, който играе ролята на посредник, позволяващ на потребителските програми да комуникират с хардуера. Тя предоставя две фундаментални свойства:
\begin{itemize}
    \item \textbf{Преносимост (Portability):} Програмите се пишат спрямо високо ниво на абстракция, за да могат да се изпълняват на различен физически хардуер без пренаписване на техния изходен код.
    \item \textbf{Делене на хардуерни ресурси:} ОС разпределя справедливо и защитено ограничените физически ресурси (памет, процесорно време, дисково пространство) между отделните програми.
\end{itemize}

Основни характеристики на съвременните ОС са \textbf{многозадачност} (multitasking --- едновременно изпълнение на няколко задачи) и \textbf{многопотребителност} (multiuser --- поддръжка на изолация между отделните потребители).

\subsection{Дефиниция на Процес}
\begin{definition}
\textbf{Процесът} е основна абстракция в операционната система, която представлява работеща (изпълняваща се) инстанция на някаква програма в реално време.
\end{definition}

За разлика от програмата, която е статична последователност от команди (файл върху диска), процесът е динамичен обект, който се характеризира със следните атрибути:
\begin{itemize}
    \item \textbf{Кой потребител го изпълнява (User ID):} Информация за собственика на процеса, определяща правата му на достъп.
    \item \textbf{Уникален идентификатор (PID):} Служебно системно число за разграничаване на процесите.
    \item \textbf{Изпълняван код:} Контекстът на инструкциите, които процесорът изпълнява.
    \item \textbf{Собствено изолирано пространство от имена (Namespace):} Включва собствена виртуална памет (стек, хип, сегмент с данни) и настройки в рамките на текущата файлова система (например таблица с дескриптори за отворени файлове).
\end{itemize}

\subsection{Комуникационни канали (IPC)}
Междупроцесната комуникация (Inter-Process Communication) е механизъм, предоставян от ОС, който позволява на два или повече процеса да обменят информация с \textbf{тяхно съгласие}.

\subsubsection{Функционални изисквания}
\begin{enumerate}
    \item \textbf{Взаимно съгласие:} Процесите трябва експлицитно да се съгласят да участват в обмена чрез използване на общи IPC обекти.
    \item \textbf{Пропускателна способност:} Ефективно пренасяне на големи количества данни с минимално натоварване (overhead) над системата.
    \item \textbf{Синхронизация:} Осигуряване на механизъм, който гарантира, че данните не се четат, преди да бъдат записани (координация във времето).
    \item \textbf{Изолация:} Неучастващи в комуникацията процеси не трябва да имат нерегламентиран достъп до обменяната информация.
\end{enumerate}

\subsubsection{Основни комуникационни абстракции}
\begin{itemize}
    \item \textbf{Споделена памет (Shared Memory):} Блок от физическата памет, който се картографира (map-ва) едновременно в адресното пространство на няколко процеса. Изключително бърз метод, тъй като се избягва копирането на данни през ядрото. Контролът и синхронизацията на достъпа (четене/писане) обаче се поемат изцяло от потребителските програми.
    \item \textbf{Комуникационна тръба (Pipe):} Прост и сигурен абстрактен механизъм за комуникация, работещ на принципа на опашката (\textbf{FIFO} --- First In, First Out). Притежава вградена синхронизация: ако тръбата е празна, четящият процес автоматично се блокира; ако се напълни, пишещият процес се приспива, докато се освободи място.
\end{itemize}

\noindent\rule{\linewidth}{0.4pt}

\section{Конкурентност и съревнование за ресурси (Race Condition)}

\subsection{Конкурентност срещу истински паралелизъм}
Начинът на изпълнение на конкурентните процеси зависи от хардуера:
\begin{itemize}
    \item \textbf{Конкурентност (Concurrency / Псевдо-паралелизъм):} Налична при еднопроцесорни системи ($1$ CPU). Диспечерът на задачите (Scheduler) бързо прекъсва и превключва контекста между процесите, създавайки илюзия за едновременна работа.
    \item \textbf{Истински паралелизъм (True Parallelism):} Наличен при многоядрени системи. Процесите се изпълняват физически едновременно върху различни процесорни ядра.
\end{itemize}

\subsection{Съревнование за ресурси (Race Condition)}
Нека имаме процесите $P$ и $Q$. Всеки от тях съдържа поредица от инструкции:
\begin{itemize}
    \item За $P$: инструкции $p_1 < p_2 < p_3$ (където $p_1$ се изпълнява хронологично преди $p_2$).
    \item За $Q$: инструкции $q_1 < q_2 < q_3$.
\end{itemize}
Понеже Диспечерът на операционната система може произволно да прекъсва изпълнението на всеки от тях, ние знаем вътрешната наредба на $p$-инструкциите и $q$-инструкциите поотделно, но \textbf{не знаем как те ще се преплетат взаимно във времето}.

\begin{definition}
\textbf{Race Condition (Съревнование за ресурси)} е ситуация, при която крайната стойност на споделен ресурс зависи от конкретния (непредвидим) ред на преплитане и изпълнение на инструкциите на конкурентните процеси от страна на Диспечера.
\end{definition}

Участъкът от програмата, в който се достъпва споделеният ресурс и където неконтролираното прекъсване е критично, се нарича \textbf{критична секция} (Critical Section).

\subsubsection{Пример за Race Condition (от лекцията на Тео)}
Нека споделената променлива е дефинирана като \lstinline|int a = 3;| в паметта.
Процесът $P$ изпълнява операция по увеличаване с $1$, а процесът $Q$ --- с $10$. На машинно ниво тези операции не са неделими (атомарни). Те се състоят от три стъпки: четене в локален регистър (\lstinline|temp|), модификация и обратно записване.

\begin{minipage}{0.45\textwidth}
\textbf{Процес P:}
\begin{lstlisting}[language=C++]
temp_P = a;
temp_P = temp_P + 1;
a = temp_P;
\end{lstlisting}
\end{minipage}
\hfill
\begin{minipage}{0.45\textwidth}
\textbf{Процес Q:}
\begin{lstlisting}[language=C++]
temp_Q = a;
temp_Q = temp_Q + 10;
a = temp_Q;
\end{lstlisting}
\end{minipage}

\vspace{0.3cm}

\textbf{Некоректен сценарий на преплитане:}
\begin{enumerate}
    \item $P$ прочита стойността на \lstinline|a|: \lstinline|temp_P = a| (стойност $3$).
    \item \textit{Диспечерът прекъсва процеса P и стартира Q!}
    \item $Q$ прочита същата стойност на \lstinline|a|: \lstinline|temp_Q = a| (стойност $3$).
    \item $Q$ модифицира стойността: \lstinline|temp_Q = 3 + 10 = 13|.
    \item $Q$ записва стойността обратно: \lstinline|a = temp_Q| (в паметта \lstinline|a| става $13$).
    \item \textit{Диспечерът прекъсва Q и връща управлението на P!}
    \item $P$ продължава от мястото на прекъсване: \lstinline|temp_P = 3 + 1 = 4|.
    \item $P$ записва стойността: \lstinline|a = temp_P| (в паметта \lstinline|a| става $4$).
\end{enumerate}
\textbf{Краен резултат:} \lstinline|a = 4|. Промяната на процес $Q$ е изцяло презаписана и загубена поради Race Condition. При коректно (последователно) изпълнение крайният резултат трябва да бъде $14$.

\noindent\rule{\linewidth}{0.4pt}

\section{Хардуерна защита на ресурс и Spinlock}

\subsection{Атомарни инструкции}
За осигуряване на взаимно изключване в критичната секция са необходими специални \textbf{атомарни} (неделими на хардуерно ниво) процесорни инструкции. Те се изпълняват в един единствен шинен такт без възможност за прекъсване. Примери са:
\begin{itemize}
    \item \lstinline|test-and-set| --- чете текущата стойност на флага и едновременно с това записва нова стойност \lstinline|true|.
    \item \lstinline|compare-and-swap| --- сравнява текущата стойност с очаквана и я променя само при съвпадение.
\end{itemize}

\subsection{Синхронизация чрез Spinlock (Въртяща се ключалка)}
\textbf{Spinlock} е нисконивоев синхронизационен примитив, при който чакащият процес се върти в непрекъснат празен цикъл (активно чакане), докато заключената променлива стане свободна.

\subsubsection{Реализация на Spinlock (Псевдокод)}
Нека имаме глобална споделена променлива: \lstinline|bool locked = false;|

\begin{lstlisting}[language=C++]
// Request to enter critical section
void lock() {
    // test_and_set atomically returns the old value of locked
    // and sets it to true. If it was true, we keep waiting.
    while (test_and_set(&locked)) {
        // busy waiting - do nothing
    }
}

// Exit critical section
void unlock() {
    locked = false;
}
\end{lstlisting}

\subsubsection{Оценка на Spinlock}
\begin{itemize}
    \item \textbf{Предимства:} Няма софтуерно закъснение за превключване на контекста (context switch) в операционната система. Изключително бърза реакция при освобождаване.
    \item \textbf{Недостатъци:} Изразходва огромно количество процесорни цикли и енергия. Напълно неприложим на едноядрени системи, защото въртящият се процес не позволява на процеса, който държи ключалката, да получи процесорно време и да я освободи.
\end{itemize}

\noindent\rule{\linewidth}{0.4pt}

\section{Синхронизация от високо ниво --- Семафор}

За да се избегне активното чакане на Spinlock при по-продължително блокирани ресурси, управлението се поема от Диспечера (Scheduler) на операционната система.

\subsection{Жизнен цикъл на процесите и Диспечер}
Диспечерът превключва състоянието на процесите между три основни фази:
\begin{enumerate}
    \item \textbf{Ready (Готов):} Процесът чака да му бъде заделено процесорно време.
    \item \textbf{Running (Изпълняващ се):} Процесът се изпълнява върху процесора в момента.
    \item \textbf{Blocked/Sleeping (Блокиран/Спящ):} Процесът чака събитие (В/И операция или освобождаване на ресурс) и не получава никакви CPU цикли.
\end{enumerate}

ОС предоставя две системни повиквания за промяна на състоянието:
\begin{itemize}
    \item \lstinline|block()| --- приспива текущия изпълняващ се процес (премества го вBlocked състояние).
    \item \lstinline|wakeup(P)| --- събужда конкретен процес $P$ (премества го в Ready състояние).
\end{itemize}

\subsection{Семафор (Semaphore)}
Семафорът е абстрактен обект за синхронизация от високо ниво, съдържащ:
\begin{itemize}
    \item Целочислен брояч (\lstinline|counter|), показващ броя на свободните единици от даден ресурс.
    \item Опашка от спящи процеси (\lstinline|queue|), чакащи за този ресурс.
\end{itemize}

\begin{definition}
\textbf{Силен семафор (Strong Semaphore)} е такъв семафор, при който опашката от спящи процеси е строго приоритетна (например FIFO опашка). Той гарантира, че няма да се стигне до системно „гладуване“ (starvation) на никой от процесите.
\end{definition}

Семафорът поддържа две атомарни операции: \lstinline|wait()| (известна като $P$) и \lstinline|signal()| (известна като $V$).

\subsubsection{Операция wait()}
Намалява стойността на брояча. Ако стойността стане отрицателна, процесът се приспива.
\begin{lstlisting}[language=C++]
void wait(Semaphore S) {
    S.counter--;
    if (S.counter < 0) {
        // Add the current process to S.queue
        block(); // Put the process to sleep
    }
}
\end{lstlisting}

\subsubsection{Операция signal()}
Увеличава стойността на брояча. Ако след увеличението стойността все още е по-малка или равна на $0$, това означава, че има спящи процеси в опашката и първият от тях се събужда.
\begin{lstlisting}[language=C++]
void signal(Semaphore S) {
    S.counter++;
    if (S.counter <= 0) {
        // Remove process P from S.queue
        wakeup(P); // Wake up process P
    }
}
\end{lstlisting}

\textit{Забележка:} Когато броячът \lstinline|S.counter| е отрицателно число (например $-k$), неговата абсолютна стойност $|S.counter| = k$ съвпада с точния брой на блокираните процеси в опашката.

\subsection{Мъртва хватка (Deadlock)}
При некоректно използване на мютекси и семафори може да се стигне до състояние на \textbf{Deadlock} (Мъртва хватка).

\begin{definition}
\textbf{Deadlock} е системна ситуация, при която два или повече процеса не могат да прогресират никога, тъй като всеки от тях е блокиран и чака освобождаването на ресурс, държан от някой друг от същите чакащи процеси.
\end{definition}

Ако имаме два процеса $P$ и $Q$ и два семафора $A = 1$, $B = 1$:
\begin{itemize}
    \item $P$ заема първо $A$, а след това чака за $B$.
    \item $Q$ заема първо $B$, а след това чака за $A$.
\end{itemize}
Ако $P$ вземе $A$, а диспечерът веднага превключи на $Q$, който вземе $B$, и двата процеса ще заспят взаимно на следващите си стъпки, чакайки ресурси, които никога няма да бъдат освободени.

\noindent\rule{\linewidth}{0.4pt}

\section{Реализация на комуникационна тръба (Pipe) с размер $n$ чрез семафори}

Това е практическо решение на класическия проблем за \textbf{Производител-Потребител (Producer-Consumer)} с ограничена опашка от $n$ елемента.

За да се позволи максимално бързо едновременно четене и писане от няколко процеса (без конфликти), в синхронизационната схема се използват два отделни мютекса:
\begin{itemize}
    \item \lstinline|mutex_write| --- защитава операцията по добавяне в опашката от други пишещи процеси.
    \item \lstinline|mutex_read| --- защитава операцията по вземане от опашката от други четящи процеси.
\end{itemize}

Останалите два семафора служат за следене на капацитета:
\begin{itemize}
    \item \lstinline|free_bytes| --- инициализира се с капацитета $n$. Намалява при добавяне на елемент.
    \item \lstinline|ready_bytes| --- инициализира се с $0$. Увеличава се при добавяне и показва колко байта са готови за четене.
\end{itemize}

\subsection{Псевдокод на решението}

\begin{lstlisting}[language=C++]
// Global objects and initialization
Semaphore free_bytes  = n; // Initially there are n free slots in the buffer
Semaphore ready_bytes = 0; // Initially there are no ready data
Semaphore mutex_write = 1; // Mutex for writing processes
Semaphore mutex_read  = 1; // Mutex for reading processes

Queue Q; // Shared queue storing up to n elements
\end{lstlisting}

\subsubsection{Процес Производител (Producer)}
\begin{lstlisting}[language=C++]
void Producer(Data item) {
    wait(free_bytes);      // Check for free slot; if none, go to sleep
    wait(mutex_write);     // Enter critical section for adding to the queue

    // --- CRITICAL SECTION ---
    Q.push(item);
    // ------------------------

    signal(mutex_write);    // Exit critical section for writing
    signal(ready_bytes);   // Signal that a new item is ready for reading
}
\end{lstlisting}

\subsubsection{Процес Потребител (Consumer)}
\begin{lstlisting}[language=C++]
Data Consumer() {
    Data item;
    wait(ready_bytes);     // Check for available data; if none, go to sleep
    wait(mutex_read);      // Enter critical section for reading from the queue

    // --- CRITICAL SECTION ---
    item = Q.pop();
    // ------------------------

    signal(mutex_read);     // Exit critical section for reading
    signal(free_bytes);    // Signal that a slot has been freed in the buffer

    return item;
}
\end{lstlisting}

\subsection{Критичен анализ на Deadlock (Мъртва хватка) при размяна на wait()}
Изключително важна особеност на решението е, че \textbf{проверката за капацитет} (\lstinline|wait(free_bytes)| / \lstinline|wait(ready_bytes)|) \textbf{задължително} трябва да се случи \textbf{преди} заемането на критичните секции (\lstinline|wait(mutex_write)| / \lstinline|wait(mutex_read)|).

Нека илюстрираме какво би се случило при некоректна размяна в кода на Производителя (при използване на общ защитен мютекс \lstinline|m| за опашката):

\begin{lstlisting}[language=C++]
// INCORRECT IMPLEMENTATION (leads to Deadlock!):
void Bad_Producer(Data item) {
    wait(m);              // Lock the queue
    wait(free_bytes);     // Check for free slot
    
    Q.push(item);
    
    signal(m);
    signal(ready_bytes);
}
\end{lstlisting}

\textbf{Сценарий за Deadlock:}
\begin{enumerate}
    \item Опашката \lstinline|Q| е пълна до своя капацитет $n$ (т.е. \lstinline|free_bytes.counter = 0|).
    \item Пристига Производител $P$. Той изпълнява успешно \lstinline|wait(m)| и заключва споделения буфер.
    \item След това $P$ извиква \lstinline|wait(free_bytes)|. Тъй като няма свободни места, броячът става $-1$ и процесът $P$ се приспива (\lstinline|block()|).
    \item В този момент пристига Потребител $C$, който иска да прочете елемент и да освободи място. За да прочете елемент обаче, той трябва първо да изпълни \lstinline|wait(m)|.
    \item Тъй като мютексът \lstinline|m| се държи от заспалия Производител $P$, Потребителят $C$ също се приспива, чакайки \lstinline|m| да се освободи.
\end{enumerate}
И двата процеса са трайно блокирани и никога няма да излязат от това състояние --- настъпил е \textbf{Deadlock}.

\end{FlushLeft}

\end{document}